\documentclass[11pt]{article}
\usepackage[margin=1.2in]{geometry}
\usepackage{amsmath,amssymb,amsthm}
\newtheorem{claim}{Claim}
\newcommand{\R}{\mathbb{R}}
\newcommand{\E}{\mathbb{E}}

\title{Tie densities are collision local times:\\ a note on the Atlas connection}
\author{Peter Cotton\thanks{Research note, \texttt{winning} repository,
\texttt{research/equilibrium/atlas\_note.tex}. Companion to the
inversion paper, whose introduction defers this connection ``for
another day.'' Citations marked [U] have not yet been verified at
source and must be before any of this is published.}}
\date{September 2026}

\begin{document}
\maketitle

\section{The static object}

The inversion paper's Proposition (photo-finish flux): for a race with
locations $\mu$ and a smooth joint density, the derivative of $i$'s win
probability with respect to a rival's location is a boundary flux,
\begin{equation}\label{eq:tie}
w_{ij} \;=\; \frac{\partial p_i}{\partial \mu_j}
\;=\; \int \varphi_{ij}(t,t)\,
\Pr\{X_k > t\ \forall k\ne i,j \mid X_i=X_j=t\}\,dt ,
\end{equation}
the probability density that $i$ and $j$ dead-heat \emph{ahead of the
field}. These are symmetric nonnegative edge conductances, and the
Jacobian is the graph Laplacian they induce.

\section{The dynamic object}

Let the performances be Brownian,
\[
X_i(t) \;=\; \mu_i + \sigma\,W_i(t),
\qquad W_1,\dots,W_n \ \text{independent},
\]
so that the time-$T$ marginals form exactly the Gaussian race of the
paper with $D_i = \sigma^2 T$: the static machinery prices
$p_i(T)=\Pr\{X_i(T)=\min_k X_k(T)\}$ and its Jacobian at every horizon,
in $O(nL)$ per horizon.

Now watch the boundary rather than the marginal. Fix a pair $(i,j)$ and
set $Y = X_i - X_j$, a Brownian motion with drift $\mu_i-\mu_j$ and
quadratic variation $\langle Y\rangle_t = 2\sigma^2 t$. The event
``$i$ and $j$ are the leading pair and tied'' is the particle system
touching $\{X_i = X_j \le X_k\ \forall k\}$. Time spent straddling that
set is measured by a local time: define the \emph{leading-pair
collision local time}
\[
\Lambda_T^{ij} \;=\;
\lim_{\varepsilon\downarrow 0}\frac{1}{2\varepsilon}
\int_0^T \mathbf 1\bigl\{|Y_s|<\varepsilon,\;
X_i(s)\wedge X_j(s) \le X_k(s)\ \forall k\ne i,j\bigr\}\,
d\langle Y\rangle_s .
\]

\begin{claim}\label{claim:main}
For the independent Brownian race,
\begin{equation}\label{eq:main}
\E\,\Lambda_T^{ij}
\;=\; 2\sigma^2 \int_0^T w_{ij}(s)\,ds
\;=\; 2\sigma^2 \int_0^T \frac{\partial p_i}{\partial\mu_j}(s)\,ds ,
\end{equation}
where $w_{ij}(s)$ is the static tie density \eqref{eq:tie} of the
time-$s$ race. Equivalently: the Jacobian entry at horizon $s$ is, up
to the factor $2\sigma^2$, the instantaneous rate at which the
leading-pair collision local time accrues.
\end{claim}

\emph{Derivation sketch.} By the occupation-times formula applied to
$Y$ under the expectation, and Fubini,
\[
\E\,\Lambda_T^{ij}
= \lim_{\varepsilon\downarrow 0}\frac{2\sigma^2}{2\varepsilon}
\int_0^T \Pr\bigl\{|Y_s|<\varepsilon,\ \text{$(i,j)$ lead}\bigr\}\,ds
= 2\sigma^2 \int_0^T \varphi^{(s)}_{Y}(0)\,
\Pr\{\text{field above}\mid \text{tie at }s\}\,ds ,
\]
and the integrand is precisely the time-$s$ tie density
\eqref{eq:tie}, i.e.\ $w_{ij}(s)$ by the paper's Proposition applied at
horizon $s$. What needs care to make this a proposition: the
$\varepsilon$-limit exchange (dominated convergence with the continuous
conditional density), and the leading-pair indicator, whose
discontinuity set is the triple-tie event --- measure zero for
nondegenerate diffusions, and exactly the degeneracy the Atlas
literature attends to [U: Ichiba--Karatzas on triple collisions].
Correlated (factor) versions keep the time-$T$ marginals inside the
covariance grammar; the identity needs only $\langle Y\rangle$ adjusted
for the factor covariance --- mechanical, unwritten. \hfill$\square$

So the boundary flux is not an \emph{analogy} to local time: for
Brownian races it \emph{is} the local-time density, and the shared
field computes it --- implicitly for all $\binom{n}{2}$ pairs, and
explicitly for any one pair --- at $O(nL)$ per horizon.

\section{Why the Atlas literature should care}

Rank-based (Atlas) models [U: Fernholz 2002; Banner--Fernholz--Karatzas
2005] assign drift and volatility by \emph{rank}; the semimartingale
decomposition of the $k$-th ranked process carries the collision local
times $\Lambda^{(k,k+1)}$ between adjacent ranks as its leakage terms,
and in stochastic portfolio theory the size effect and the attrition of
the largest market weight are expressed through them. Market weights
are a choice-probability vector --- the inversion paper's
portfolio remark is the same bridge --- and rank-crossing local times
are what the static machinery evaluates, one horizon at a time.

Trade flows in both directions.

\begin{enumerate}
\item \emph{Computation to Atlas.} Collision local times are
analytically available only for small $n$ or special symmetric cases
[U]. Claim~\ref{claim:main} plus the lattice evaluates expected
local-time densities for the leading pair at $O(nL)$ per horizon for
$n$ in the millions; the removal counterfactual (delete the leaders,
re-race) reaches lower ranks. Checkable against simulated ranked
Brownian systems --- an exact-versus-Monte-Carlo experiment of the kind
the repository already referees.
\item \emph{Structure to us.} The Atlas decomposition names the dynamic
inversion: recover drifts from observed \emph{leadership turnover}
(occupation and crossing rates of ranks), of which the static inversion
is the one-horizon shadow.
\end{enumerate}

\section{Scope honesty}

The identity is stated for the leading pair; adjacent-rank local times
at rank $k$ need a top-$k$ cavity (delete the $k-1$ leaders and
re-race), with cost and accuracy unmeasured. Nothing here addresses
stationary gap distributions [U: Pal--Pitman], where much of the Atlas
literature lives; the contribution in prospect is finite-horizon and
computational. Every bracketed citation above is unverified at source.

\end{document}
